






\documentclass[sigconf]{aamas} 


\usepackage{thm-restate}


\newcommand{\epsdeltaone}{\epsilon_1}
\newcommand{\bestparetopayoffi}{u^*_i}
\newcommand{\bestparetopayoffj}{u^*_j}
\newcommand{\bestparetopayoffone}{u^*_1}
\newcommand{\bestparetopayofftwo}{u^*_2}
\newcommand{\fullresp}{R}
\newcommand{\history}{h}
\newcommand{\ponefav}{a_1}
\newcommand{\ptwofav}{a_2}
\newcommand{\cost}{\alpha}
\newcommand{\prop}{p}
\newcommand{\prior}{q}
\newcommand{\demandset}{\mathcal{D}}
\newcommand{\reasonset}{\mathcal{R}}
\newcommand{\alloc}{x}
\newcommand{\allocspace}{\mathcal{X}}
\newcommand{\impartial}{\mathcal{I}}
\newcommand{\newactspacetwo}{\tilde{\mathcal{A}}_2}
\newcommand{\allprogrn}{\allprog^{\text{rn}}}
\newcommand{\allprogreneg}{\allprog^{\text{RN}}}
\newcommand{\prnsup}{\textnormal{PRN}}
\newcommand{\prnspacei}{\mathcal{P}_i}
\newcommand{\prnspacej}{\mathcal{P}_j}
\newcommand{\prnspaceone}{\mathcal{P}_1}
\newcommand{\prnspacetwo}{\mathcal{P}_2}
\newcommand{\allprogprn}{\allprog^{\prnsup}}
\newcommand{\defrenegi}{{\renegi}^{\textnormal{def}}}
\newcommand{\defrenegj}{{\renegj}^{\textnormal{def}}}
\newcommand{\defrenegone}{{\renegone}^{\textnormal{def}}}
\newcommand{\defrenegtwo}{{\renegtwo}^{\textnormal{def}}}
\newcommand{\dumbset}{Z}
\newcommand{\disone}{d_1}
\newcommand{\distwo}{d_2}
\newcommand{\disi}{d_i}
\newcommand{\disj}{d_j}
\newcommand{\demi}{x_i}
\newcommand{\demj}{x_j}
\newcommand{\demone}{x_1}
\newcommand{\demtwo}{x_2}
\newcommand{\respi}{r_i}
\newcommand{\respj}{r_j}
\newcommand{\resp}{r}
\newcommand{\probfail}{p}
\newcommand{\costthreshold}{\underline{\alpha}}
\newcommand{\fulldisc}{Y}
\newcommand{\devspace}{P}
\newcommand{\indicator}[1]{\mathbb{I}[#1]}
\newcommand{\thr}{t}
\newcommand{\belief}{Q}
\newcommand{\tademand}{Y}
\newcommand{\baseta}{y}

\newcommand{\selection}{sel}%{bar}
\newcommand{\renegb}{\mathbf{\selection}}




\newcommand{\dirac}[1]{\delta_{#1}}
\newcommand{\observed}{T}
\newcommand{\progpayiparetomin}{\progpayi^{\vee}}

\usepackage{amsthm}

\bibliographystyle{ACM-Reference-Format}

\theoremstyle{definition}
\newtheorem{property}{Property}
\newtheorem{definition}{Definition} % Create the "definition" environment
\newtheorem{assumption}[definition]{Assumption}
\setcounter{property}{0}

\newtheorem*{assum:nopunish}{Assumption 7}



\newcommand{\mult}{1.15}
\newcommand{\offs}{0.05}
\newcommand{\globalflag}{0}
\newcommand{\pointx}{0.8}%0.65}
\newcommand{\tritop}{2.4}%1.95} % 3 x pointx
\newcommand{\tribottom}{0.267}%0.22} % pointx / 3

% --- begin inlined defaults.sty ---
\usepackage{graphicx,psfrag,amsmath,amsfonts,verbatim}
\usepackage{amsmath,amsfonts}%,
\usepackage{changepage}
\usepackage{latexsym, amscd, amsfonts, eucal, mathrsfs, amsmath, %amssymb,
amsthm, %xypic,
xr, stmaryrd, color, enumerate, mathtools}
\usepackage[utf8]{inputenc}
\newcommand\myeq{\stackrel{\mathclap{\normalfont\mbox{a.s.}}}{=}}
\newcommand{\expectation}[4]{\mathbb{E}_{#1 \sim #2(\cdot | #3)}#4}
\newcommand{\dist}{\mathcal{D}}
\usepackage{xcolor}
\definecolor{lightergray}{gray}{0.9}

\usepackage{natbib}

\newcommand{\multilinecomment}[1]{}

\DeclareSymbolFont{bbold}{U}{bbold}{m}{n}
\DeclareSymbolFontAlphabet{\mathbbold}{bbold}

\DeclareMathOperator*{\argmin}{arg\,min}
\DeclareMathOperator*{\argmax}{arg\,max}
\DeclarePairedDelimiter{\ceil}{\lceil}{\rceil}
\DeclarePairedDelimiter{\floor}{\lfloor}{\rfloor}


\newcommand{\littler}{rn}%{r}
\newcommand{\pmsup}{\textnormal{PMM}}%\text{PMM}
\newcommand{\svsup}{\textnormal{RN}}
\newcommand{\itsup}{\textnormal{IRN}}
\newcommand{\csrsup}{\textnormal{RN}}
\newcommand{\modcsrsup}{\widetilde{\textnormal{CS}}}
\newcommand{\rensup}{\textnormal{rn}}
\newcommand{\intsect}{\textbf{AS}}%C}
\newcommand{\defprog}[1]{{#1}^{\textnormal{def}}}
\newcommand{\defprogi}{\progi^{\textnormal{def}}}
\newcommand{\defprogone}{\progone^{\textnormal{def}}}
\newcommand{\defprogtwo}{\progtwo^{\textnormal{def}}}
\newcommand{\defprogj}{\progj^{\textnormal{def}}}\newcommand{\defprogmi}{\deffullprog_{-i}}
\newcommand{\defprogmj}{\deffullprog_{-j}}
\newcommand{\defprogmji}{\progmj^{\defprogi}}%{\defprog{\progmj}^{i}}
\newcommand{\defeqprogmji}{\progmj^{\defprog{\irfsprogi(\eqprogi)}}}
\newcommand{\renegbi}{\mathbf{\selection}^i}
\newcommand{\renegbmi}{\mathbf{\selection}^{-i}}
\newcommand{\renegbj}{\mathbf{\selection}^j}
\newcommand{\bargdefault}{\mathbf{ro}}
\newcommand{\fullrenegi}{\renegi({K})}%{\underline{\renegi}}
\newcommand{\fullrenegj}{\renegj({K})}%{\underline{\renegj}}
\newcommand{\pmp}{\mathbf{PMP}}%{M}
\newcommand{\yudproj}{\pmp}
\newcommand{\yudproji}{\pmp_i}
\newcommand{\yudprojj}{\pmp_j}
\newcommand{\yudprojone}{\pmp_1}
\newcommand{\yudprojtwo}{\pmp_2}
\newcommand{\termemph}[1]{\textbf{#1}}%{\textit{#1}}

\newcommand{\fsspaceireni}{\fsspacei(\renegi)}%{\fsspacei(\renegi, \renegb)}
\newcommand{\fsspaceinewreni}{\fsspacei(\newrenegi)}%{\fsspacei(\newrenegi, \renegb)}
\newcommand{\itspaceireni}{\itspacei(\fullrenegi)}%{\itspacei(\fullrenegi, \renegb)}
\newcommand{\fullfsspacei}{\fsspacei}%{\fsspacei(\renegb)}
\newcommand{\fullfsspaceone}{\fsspaceone}%{\fsspacei(\renegb)}
\newcommand{\fullfsspacetwo}{\fsspacetwo}%{\fsspacei(\renegb)}
\newcommand{\fullitspacei}{\itspacei}%{\itspacei(\renegb)}
\newcommand{\genfunc}{s}

\newcommand{\corrpolicy}{\boldsymbol{\tau}}
\newcommand{\gencorr}{\boldsymbol{\mu}}
\newcommand{\jointstrategy}{\boldsymbol{a}}
\newcommand{\fulltype}{\mathbf{t}}
\newcommand{\randomization}{\mathcal{C}}
\newcommand{\randdomain}{R}
\newcommand{\punishment}{\boldsymbol{\underline{\tau}}}
\newcommand{\wrongguess}{y}
\newcommand{\exploited}{x}
\newcommand{\tildeps}{\tilde \epsilon}
\newcommand{\tildtau}{\tilde \tau}
\newcommand{\overB}{\overline{B}}
\newcommand{\underB}{\underline{B}}
\newcommand{\genci}{\tilde{c}_i}
\newcommand{\gencj}{\tilde{c}_j}
\newcommand{\genpi}{\widetilde{p}_i}
\newcommand{\genpj}{\widetilde{p}_j}
\newcommand{\gencone}{\tilde{c}_1}
\newcommand{\genctwo}{\tilde{c}_2}
\newcommand{\genpone}{\tilde{p}_1}
\newcommand{\genptwo}{\tilde{p}_2}
\newcommand{\think}{t}
\newcommand{\actprogone}[1]{\act_1^{#1}}
\newcommand{\actprogonest}[1]{\act_1^{*#1}}
\newcommand{\actprogtwo}[1]{\act_2^{#1}}
\newcommand{\actprogtwost}[1]{\act_2^{*#1}}
\newcommand{\tiebreakone}{{\text{BR}}_1}
\newcommand{\tiebreaktwo}{{\text{BR}}_2}
\newcommand{\tiebreaki}{{\text{BR}}_i}
\newcommand{\tiebreakj}{{\text{BR}}_j}
\newcommand{\gentiebreakone}{\widetilde{{\text{BR}}}_1}
\newcommand{\gentiebreaktwo}{\widetilde{{\text{BR}}}_2}
\newcommand{\gentiebreaki}{\widetilde{{\text{BR}}}_i}
\newcommand{\gentiebreakj}{\widetilde{{\text{BR}}}_j}
\newcommand{\tieseti}{\mathcal{B}_i}
\newcommand{\tiesetj}{\mathcal{B}_j}
\newcommand{\fulltie}{\textbf{{\text{BR}}}}
\newcommand{\genfullact}{\tilde{\fullact}}
\newcommand{\expset}{\mathcal{T}}
\newcommand{\decproci}{D_i}
\newcommand{\decprocj}{D_j}
\newcommand{\decprocone}{D_1}
\newcommand{\decproctwo}{D_2}
\newcommand{\fulldecproc}{\textbf{D}}
\newcommand{\decspace}{\mathcal{D}}
\newcommand{\decspaceone}{\mathcal{D}_1}
\newcommand{\decspacetwo}{\mathcal{D}_2}
\newcommand{\decspacei}{\mathcal{D}_i}
\newcommand{\decspacej}{\mathcal{D}_j}
\newcommand{\yudset}{Y}
\newcommand{\compyudset}{\tilde{Y}}
\newcommand{\fullrenegspacej}{\mathcal{r}^{j}}
\newcommand{\svrenegspace}{\mathcal{R}}
\newcommand{\svrenegspacei}{\mathcal{R}^{i}}
\newcommand{\svrenegspacej}{\mathcal{R}^{j}}
\newcommand{\fullsvspacei}{\underline{\svrenegspacei}}
\newcommand{\fullrenegspace}{\mathcal{R}}
\newcommand{\fullsvspace}{\mathcal{R}^{\svsup}}
\newcommand{\fullsvspacej}{\underline{\svrenegspacej}}
\newcommand{\uniqyudproj}{y}
\newcommand{\overyud}{\overline{Y}}
\newcommand{\deffullact}{\widehat{\fullact}}
\newcommand{\defacti}{\widehat{a}_i}
\newcommand{\defactj}{\widehat{a}_j}
\newcommand{\defactmi}{\deffullact_{-i}}
\newcommand{\acceptone}{S_1}
\newcommand{\accepttwo}{S_2}
\newcommand{\accepti}{S_i}
\newcommand{\acceptj}{S_j}
\newcommand{\renega}{\reneg^{\text{I}}}
\newcommand{\renegk}{\reneg^{(k)}}
\newcommand{\renegkp}{\reneg^{(k+1)}}
\newcommand{\renegai}{\reneg^{i,\text{I}}}
\newcommand{\renegaj}{\reneg^{j,\text{I}}}
\newcommand{\renegki}{\reneg^{i,(k)}}
\newcommand{\renegkim}{\reneg^{-i,(k)}}
\newcommand{\renegkjm}{\reneg^{-j,(k)}}
\newcommand{\renegkj}{\reneg^{j,(k)}}
\newcommand{\renegonei}{\reneg^{i,(1)}}
\newcommand{\renegonej}{\reneg^{j,(1)}}
\newcommand{\renegbki}{\reneg^{i,(K)}}
\newcommand{\renegbkj}{\reneg^{j,(K)}}
\newcommand{\renegtwoi}{\reneg^{i,(2)}}
\newcommand{\renegtwoj}{\reneg^{j,(2)}}
\newcommand{\renegkpi}{\reneg^{i,(k+1)}}
\newcommand{\renegkpj}{\reneg^{j,(k+1)}}
\newcommand{\renegkmi}{\reneg^{i,(k-1)}}
\newcommand{\renegkmj}{\reneg^{j,(k-1)}}
\newcommand{\renegkmmi}{\reneg^{i,(k-2)}}
\newcommand{\renegkmmj}{\reneg^{j,(k-2)}}
\newcommand{\renegaone}{\reneg^{1,\text{I}}}
\newcommand{\renegatwo}{\reneg^{2,\text{I}}}
\newcommand{\renegkone}{\reneg^{1,(k)}}
\newcommand{\renegktwo}{\reneg^{2,(k)}}
\newcommand{\renegoneone}{\reneg^{1,(1)}}
\newcommand{\renegonetwo}{\reneg^{2,(1)}}
\newcommand{\renegbkone}{\reneg^{1,(K)}}
\newcommand{\renegbktwo}{\reneg^{2,(K)}}
\newcommand{\renegtwoone}{\reneg^{1,(2)}}
\newcommand{\renegtwotwo}{\reneg^{2,(2)}}
\newcommand{\renegkpone}{\reneg^{1,(k+1)}}
\newcommand{\renegkptwo}{\reneg^{2,(k+1)}}
\newcommand{\renegkmone}{\reneg^{1,(k-1)}}
\newcommand{\renegkmtwo}{\reneg^{2,(k-1)}}
\newcommand{\renegc}{\reneg^{\text{II}}}
\newcommand{\renegci}{\reneg^{i,\text{II}}}
\newcommand{\renegcj}{\reneg^{j,\text{II}}}
\newcommand{\newrenega}{{\widetilde{\reneg}}^{\text{I}}}
\newcommand{\pmm}{\fullpay^{\pmsup}}
\newcommand{\actpmm}{\fullact^{\pmsup}}
\newcommand{\pmmi}{\payi^{\pmsup}}
\newcommand{\pmmip}{\payip^{\pmsup}}
\newcommand{\pmmj}{\payj^{\pmsup}}
\newcommand{\pmmone}{\payone^{\pmsup}}
\newcommand{\pmmtwo}{\paytwo^{\pmsup}}
\newcommand{\pmmn}{u_n^{\pmsup}}
\newcommand{\klist}{\{1,\dots,{K}\}}
\newcommand{\kmlist}{\{1,\dots,{K}-1\}}
\newcommand{\newrenegb}{\mathbf{\widetilde{B}}}
\newcommand{\newrenegc}{{\widetilde{\reneg}}^{\text{II}}}
\newcommand{\newreneg}{{\widetilde{\reneg}}}
\newcommand{\newrenegi}{{\widetilde{\reneg}}^{i}}
\newcommand{\newrenegimj}{\renegmj_{\newrenegi}}
\newcommand{\newrenegj}{{\widetilde{\reneg}}^{j}}
\newcommand{\newrenegmi}{{\widetilde{\reneg}}^{-i}}
\newcommand{\newrenegia}{{\widetilde{\reneg}}^{i,\text{{I}}}}
\newcommand{\newrenegki}{{\widetilde{\reneg}}^{i,(k)}}
\newcommand{\newrenegkj}{{\widetilde{\reneg}}^{j,(k)}}
\newcommand{\newrenegone}{{\widetilde{\reneg}}^{1}}
\newcommand{\newrenegonei}{{\widetilde{\reneg}}^{i,(1)}}
\newcommand{\newrenegonej}{{\widetilde{\reneg}}^{j,(1)}}
\newcommand{\newrenegoneone}{{\widetilde{\reneg}}^{1,(1)}}
\newcommand{\newrenegbki}{{\widetilde{\reneg}}^{i,(K)}}
\newcommand{\newrenegbkj}{{\widetilde{\reneg}}^{j,(K)}}
\newcommand{\newrenegbkone}{{\widetilde{\reneg}}^{1,(K)}}
\newcommand{\newrenegtwoi}{{\widetilde{\reneg}}^{i,(2)}}
\newcommand{\newrenegtwoj}{{\widetilde{\reneg}}^{j,(2)}}
\newcommand{\newrenegkpi}{{\widetilde{\reneg}}^{i,(k+1)}}
\newcommand{\newrenegkpj}{{\widetilde{\reneg}}^{j,(k+1)}}
\newcommand{\newrenegkmi}{{\widetilde{\reneg}}^{i,(k-1)}}
\newcommand{\newrenegkmj}{{\widetilde{\reneg}}^{j,(k-1)}}
\newcommand{\newrenegic}{{\widetilde{\reneg}}^{i,\text{{II}}}}
\newcommand{\dominate}{D}
\newcommand{\worsei}{W_i}
\newcommand{\worsej}{W_j}
\newcommand{\deffullprog}{\defprog{\fullprog}}
\newcommand{\defnewfullprog}{{\newfullprog}^{\textnormal{def}}}
\newcommand{\defnewfullprogi}{{\widetilde{\textbf{\prog}}^{i^{\textnormal{def}}}}}
\newcommand{\renegx}{\reneg^{X}}
\newcommand{\renegi}{\reneg^{i}}
\newcommand{\renegj}{\reneg^{j}}
\newcommand{\renegmi}{\reneg^{-i}}
\newcommand{\renegmj}{\reneg^{-j}}
\newcommand{\renegone}{\reneg^{1}}
\newcommand{\renegtwo}{\reneg^{2}}
\newcommand{\genreneg}{\renegj}
\newcommand{\genrenega}{\renegaj}
\newcommand{\genrenegc}{\renegcj}
\newcommand{\genrenegx}{\overline{\reneg}^{X}}
\newcommand{\renegsta}{{\renegst}^{\text{I}}}
\newcommand{\renegstb}{\renegb^*}
\newcommand{\renegstc}{{\renegst}^{\text{II}}}
\newcommand{\renegstx}{{\renegst}^{X}}
\newcommand{\renai}{\reni^{\text{I}}}
\newcommand{\renbi}{B_i}
\newcommand{\renci}{\reni^{\text{II}}}
\newcommand{\rencii}{\reni^{i,\text{II}}}
\newcommand{\renii}{\reni^{i}}
\newcommand{\renaii}{\reni^{i,\text{I}}}
\newcommand{\renxi}{\reni^{X}}
\newcommand{\renaj}{\renj^{\text{I}}}
\newcommand{\renbj}{B_j}
\newcommand{\rencj}{\renj^{\text{II}}}
\newcommand{\renxj}{\renj^{X}}
\newcommand{\powerset}{\mathbf{C}}
\newcommand{\effic}{{E}}
\newcommand{\fullx}{\fullact^{x}}
\newcommand{\fully}{\fullact^{y}}
\newcommand{\fullz}{\fullact^{z}}
\newcommand{\fullw}{\fullact^{w}}
\newcommand{\game}{G}

\definecolor{darkgreen}{rgb}{0,0.5,0}
\newcommand{\kibitz}[2]{\ifnum\Comments=1\textcolor{#1}{#2}\fi}
\newcommand{\adigi}[1]{\kibitz{blue}{[AD: #1]}}
\newcommand{\jc}[1]{\kibitz{red}{[JC: #1]}}

\newcommand{\bctransform}[1]{{#1}^B}

\newcommand{\commentedsection}[2]{%
    \ifnum#1=1
        \ignorespaces % Ignores any leading spaces
        #2
        \ignorespacesafterend % Ignores any trailing spaces
    \else
        \ignorespaces
        \ignorespacesafterend
    \fi
}

\newcommand{\selfunction}[0]{selection}%{bargaining}
\newcommand{\notsafe}{N}
\newcommand{\maximini}{u_i^M}
\newcommand{\maximinj}{u_j^M}
\newcommand{\maximinip}{u_{i'}^M}
\newcommand{\maximinone}{u_1^M}
\newcommand{\maximintwo}{u_2^M}
\newcommand{\procspace}{\Pi}
\newcommand{\actdprogone}{\act^{\defprog{\progone}}}
\newcommand{\typeinformation}{\Theta}
\newcommand{\minusi}{_{-i}}
\newcommand{\minusj}{_{-j}}
\newcommand{\programspace}{P}
\newcommand{\progmone}{\fullprog_{-1}}
\newcommand{\progi}{\prog_i}
\newcommand{\progj}{\prog_j}
\newcommand{\progmi}{\fullprog_{-i}}
\newcommand{\progmj}{\fullprog_{-j}}
\newcommand{\ogactone}{\underline{a}_1}
\newcommand{\actmisone}{\actspaceone^{\mathcal{M}}}
\newcommand{\actmistwo}{\actspacetwo^{\mathcal{M}}}
\newcommand{\actmis}{\actspace^{\mathcal{M}}}
\newcommand{\progmisi}{\programspace_i^{\mathcal{M}}}
\newcommand{\progmisj}{\programspace_j^{\mathcal{M}}}
\newcommand{\progmisone}{\programspace_1^{\mathcal{M}}}
\newcommand{\progmistwo}{\programspace_2^{\mathcal{M}}}
\newcommand{\allprog}{\mathcal{P}}
\newcommand{\fsallprog}[1]{\mathcal{P}^{#1}}
\newcommand{\fsfullprior}[1]{\fullprior^{#1}}
\newcommand{\progspacei}{\programspace_i}
\newcommand{\progspacej}{\programspace_j}
\newcommand{\progspaceone}{\programspace_1}
\newcommand{\progspacetwo}{\programspace_2}
\newcommand{\progspacen}{\programspace_n}
\newcommand{\bprogmisi}{\bctransform{\programspace}_i^{\mathcal{M}}}
\newcommand{\bprogmisj}{\bctransform{\programspace}_j^{\mathcal{M}}}
\newcommand{\bprogmisone}{\bctransform{\programspace}_1^{\mathcal{M}}}
\newcommand{\bprogmistwo}{\bctransform{\programspace}_2^{\mathcal{M}}}
\newcommand{\ballprog}{\bctransform{\mathcal{\programspace}}}
\newcommand{\bprogspacei}{\bctransform{\programspace}_i}
\newcommand{\bprogspacej}{\bctransform{\programspace}_j}
\newcommand{\bprogspaceone}{\bctransform{\programspace}_1}
\newcommand{\bprogspacetwo}{\bctransform{\programspace}_2}
\newcommand{\progbrone}{\prog_1^{BR}}
\newcommand{\progbrtwo}{\prog_2^{BR}}
\newcommand{\progbri}{\prog_i^{BR}}
\newcommand{\progbrj}{\prog_j^{BR}}
\newcommand{\cespaceone}{\programspace_1^{CE}}
\newcommand{\cespacetwo}{\programspace_2^{CE}}
\newcommand{\actoneone}{\actone^1}
\newcommand{\actonetwo}{\actone^2}
\newcommand{\poltwo}{\pi_2}
\newcommand{\fsspaceone}{\programspace_1^{\csrsup}}
\newcommand{\fsspacetwo}{\programspace_2^{\csrsup}}
\newcommand{\onefsspacei}{\programspace_i^{\rensup}}
\newcommand{\onefsspaceip}{\programspace_{i'}^{\rensup}}
\newcommand{\onefsspacej}{\programspace_j^{\rensup}}
\newcommand{\onefsspaceone}{\programspace_1^{\rensup}}
\newcommand{\onefsspacetwo}{\programspace_2^{\rensup}}
\newcommand{\fullrenegone}{\underline{\renegone}}\newcommand{\fullrenegtwo}{\underline{\renegtwo}}
\newcommand{\fsspacei}{\programspace_i^{\csrsup}}
\newcommand{\fsspacej}{\programspace_j^{\csrsup}}
\newcommand{\modspacei}{\programspace_i^{\modcsrsup}}
\newcommand{\modspacej}{\programspace_j^{\modcsrsup}}
\newcommand{\itspacei}{\programspace_i^{\itsup}}
\newcommand{\itspacej}{\programspace_j^{\itsup}}
\newcommand{\itspaceone}{\programspace_1^{\itsup}}
\newcommand{\itspacetwo}{\programspace_2^{\itsup}}
\newcommand{\ufsspaceone}{\programspace_1^{E}}
\newcommand{\ufsspacetwo}{\programspace_2^{E}}
\newcommand{\ufsspacei}{\programspace_i^{E}}
\newcommand{\ufsspacej}{\programspace_j^{E}}
\newcommand{\unispaceone}{\programspace_1^{U}}
\newcommand{\unispacetwo}{\programspace_2^{U}}
\newcommand{\unispacei}{\programspace_i^{U}}
\newcommand{\unispacej}{\programspace_j^{U}}
\newcommand{\spispaceone}{\programspace_1^{S}}
\newcommand{\spispacetwo}{\programspace_2^{S}}
\newcommand{\spispacei}{\programspace_i^{S}}
\newcommand{\spispacej}{\programspace_j^{S}}
\newcommand{\nash}[1]{\text{NE}(#1)}
\newcommand{\subjeq}{\text{SE}}
\newcommand{\notfsspacei}{\overline{\programspace_i^{\rensup}}}
\newcommand{\notfsspacej}{\overline{\programspace_j^{\rensup}}}
\newcommand{\notfsspaceone}{\overline{\programspace_1^{\rensup}}}
\newcommand{\notfsspacetwo}{\overline{\programspace_2^{\rensup}}}
\newcommand{\bcrspaceone}{\programspace_1^{D}}
\newcommand{\bcrspacetwo}{\programspace_2^{D}}
\newcommand{\fullprogpay}{\textbf{U}}
\newcommand{\fullpay}{\textbf{u}}
\newcommand{\bcrspacei}{\programspace_i^{D}}
\newcommand{\bcrspacej}{\programspace_j^{D}}
\newcommand{\neweqprogi}{\widetilde{\eqprogi}}
\newcommand{\neweqprogj}{\widetilde{\eqprogj}}
\newcommand{\neweqprogone}{\widetilde{\eqprogone}}
\newcommand{\neweqprogtwo}{\widetilde{\eqprogtwo}}
\newcommand{\newprogi}{\widetilde{\progi}}
\newcommand{\newprogj}{\widetilde{\progj}}
\newcommand{\newprogone}{\widetilde{\progone}}
\newcommand{\newprogtwo}{\widetilde{\progtwo}}
\newcommand{\procedure}{\pi}
\newcommand{\reneg}{\mathbf{RN}}
\newcommand{\smreneg}{\mathbf{\littler}}
\newcommand{\smrenegi}{\mathbf{\littler}^i}
\newcommand{\smrenegj}{\mathbf{\littler}^j}
\newcommand{\smrenegone}{\mathbf{\littler}^1}
\newcommand{\smrenegtwo}{\mathbf{\littler}^2}
\newcommand{\renegst}{\reneg^*}
\newcommand{\renone}{\textnormal{\littler}_1}
\newcommand{\rentwo}{\textnormal{\littler}_2}
\newcommand{\reni}{\textnormal{\littler}_i}
\newcommand{\renj}{\textnormal{\littler}_j}
\newcommand{\renstone}{R^*_1}
\newcommand{\rensttwo}{R^*_2}
\newcommand{\rensti}{R^*_i}
\newcommand{\renstj}{R^*_j}
\newcommand{\fsprogone}{f_1}
\newcommand{\fsprogtwo}{f_2}
\newcommand{\fsprogi}{f_i}
\newcommand{\fsprogj}{f_j}
\newcommand{\irfsprogone}{f_1^*}
\newcommand{\irfsprogtwo}{f_2^*}
\newcommand{\irfsprogi}{f_i^*}
\newcommand{\irfsprogj}{f_j^*}
\newcommand{\fullact}{\jointstrategy}
\newcommand{\fullactst}{\jointstrategy^*}
\newcommand{\fullactunder}{\underline{\jointstrategy}}
\newcommand{\fullactover}{\overline{\jointstrategy}}
\newcommand{\actiover}{\overline{\acti}}
\newcommand{\actjover}{\overline{\actj}}
\newcommand{\actoneover}{\overline{\actone}}
\newcommand{\acttwoover}{\overline{\acttwo}}
\newcommand{\actist}{\act_i^*}
\newcommand{\actjst}{\act_j^*}
\newcommand{\actiunder}{\underline{\act}_i}
\newcommand{\actjunder}{\underline{\act}_j}
\newcommand{\actip}{\act_{i'}}
\newcommand{\actjp}{\act_{j'}}
\newcommand{\actonest}{\act_1^*}
\newcommand{\acttwost}{\act_2^*}
\newcommand{\bfsspaceone}{\bctransform{\programspace}_1^{F}}
\newcommand{\bfsspacetwo}{\bctransform{\programspace}_2^{F}}
\newcommand{\progiunder}{\underline{\prog}_i}
\newcommand{\progjunder}{\underline{\prog}_j}
\newcommand{\progoneunder}{\underline{\prog}_1}
\newcommand{\progtwounder}{\underline{\prog}_2}
\newcommand{\progipunder}{\underline{\prog}_{i'}}
\newcommand{\progjpunder}{\underline{\prog}_{j'}}
\newcommand{\simulf}{s_F}
\newcommand{\simulb}{s_B}
\newcommand{\simulr}{s_R}
\newcommand{\follower}{F}
\newcommand{\actbrtwo}{\acttwo^{BR}}
\newcommand{\fulltheta}{\boldsymbol{\theta}}
\newcommand{\fullbehprob}{\textbf{\prior}^B}
\newcommand{\parti}{F_i}
\newcommand{\partj}{F_j}
\newcommand{\partone}{F_1}
\newcommand{\parttwo}{F_2}
\newcommand{\behprobi}{\prior_i^B}
\newcommand{\behprobj}{\prior_j^B}
\newcommand{\behprobone}{\prior_1^B}
\newcommand{\behprobtwo}{\prior_2^B}
\newcommand{\ponece}{\progone^{CE}}
\newcommand{\ptwoce}{\progtwo^{CE}}
\newcommand{\bonece}{\burnone^{CE}}
\newcommand{\btwoce}{\burntwo^{CE}}
\newcommand{\eqprogone}{\prog_1^*}
\newcommand{\eqprogtwo}{\prog_2^*}
\newcommand{\eqprogn}{\prog_n^*}
\newcommand{\acteqone}{\act_1^*}
\newcommand{\acteqtwo}{\act_2^*}
\newcommand{\acteqi}{\act_i^*}
\newcommand{\acteqj}{\act_j^*}
\newcommand{\fullfs}{\textbf{f}}
\newcommand{\irfullfs}{\textbf{f}^*}
\newcommand{\eqprogi}{\prog_i^*}
\newcommand{\eqprogip}{\prog_{i'}^*}
\newcommand{\eqprogj}{\prog_j^*}
\newcommand{\comspacei}{C_i}
\newcommand{\comspacej}{C_j}
\newcommand{\comspaceone}{C_1}
\newcommand{\comspacetwo}{C_2}
\newcommand{\burnspacei}{B_i}
\newcommand{\burnspacej}{B_j}
\newcommand{\burnspaceone}{B_1}
\newcommand{\burnspacetwo}{B_2}
\newcommand{\thetab}{\theta^B}
\newcommand{\thetar}{\theta^R}
\newcommand{\thetaf}{\theta^F}
\newcommand{\thetai}{\theta_i}
\newcommand{\thetaj}{\theta_j}
\newcommand{\thetaone}{\theta_1}
\newcommand{\thetatwo}{\theta_2}
\newcommand{\priordomaini}{\mathcal{Q}_i}
\newcommand{\priordomainj}{\mathcal{Q}_j}
\newcommand{\priordomainij}{\mathcal{Q}_{ij}}
\newcommand{\priordomainji}{\mathcal{Q}_{ji}}
\newcommand{\bpriordomaini}{\bctransform{\mathcal{Q}}_i}
\newcommand{\bpriordomainj}{\bctransform{\mathcal{Q}}_j}
\newcommand{\bpriordomainij}{\bctransform{\mathcal{Q}}_{ij}}
\newcommand{\bpriordomainji}{\bctransform{\mathcal{Q}}_{ji}}
\newcommand{\bpostdomainij}{\bctransform{\mathfrak{Q}}_{ij}}
\newcommand{\bpostdomainji}{\bctransform{\mathfrak{Q}}_{ji}}
\newcommand{\bpostdomainonetwo}{\bctransform{\mathfrak{Q}}_{12}}
\newcommand{\bpostdomaintwoone}{\bctransform{\mathfrak{Q}}_{21}}
\newcommand{\mpriordomainij}{\bctransform{\mathcal{M}}_{ij}}
\newcommand{\mpriordomainji}{\bctransform{\mathcal{M}}_{ji}}
\newcommand{\mpriordomainonetwo}{\bctransform{\mathcal{M}}_{12}}
\newcommand{\mpriordomaintwoone}{\bctransform{\mathcal{M}}_{21}}
\newcommand{\priordomainonetwo}{\mathcal{Q}_{12}}
\newcommand{\priordomaintwoone}{\mathcal{Q}_{21}}
\newcommand{\bpriordomainonetwo}{\bctransform{\mathcal{Q}}_{12}}
\newcommand{\bpriordomaintwoone}{\bctransform{\mathcal{Q}}_{21}}
\newcommand{\bbeldomainij}{\widetilde{\bpriordomainij}}
\newcommand{\bbeldomainji}{\widetilde{\bpriordomainji}}
\newcommand{\bbeldomainonetwo}{\widetilde{\bpriordomainonetwo}}
\newcommand{\bbeldomaintwoone}{\widetilde{\bpriordomaintwoone}}
\newcommand{\actspace}{A}
\newcommand{\allact}{\mathcal{A}}
\newcommand{\actspacei}{A_i}
\newcommand{\actspacej}{A_j}
\newcommand{\actspaceone}{A_1}
\newcommand{\actspacetwo}{A_2}
\newcommand{\actspacen}{A_n}
\newcommand{\commiti}{c_i}
\newcommand{\commitj}{c_j}
\newcommand{\commitij}{c_i^j}
\newcommand{\commitjj}{c_j^j}
\newcommand{\commitii}{c_i^i}
\newcommand{\commitji}{c_j^i}
\newcommand{\commitone}{c_1}
\newcommand{\committwo}{c_2}
\newcommand{\commitonetwo}{c_1^2}
\newcommand{\committwotwo}{c_2^2}
\newcommand{\commitoneone}{c_1^1}
\newcommand{\committwoone}{c_2^1}
\newcommand{\burni}{b_i}
\newcommand{\burnj}{b_j}
\newcommand{\burnone}{b_1}
\newcommand{\burntwo}{b_2}
\newcommand{\priori}{\beta_i}%{q_i}
\newcommand{\priorj}{\beta_j}%{q_j}
\newcommand{\genpriorij}{\tilde{q}_{ij}}
\newcommand{\genpriorji}{\tilde{q}_{ji}}
\newcommand{\genprioronetwo}{\tilde{q}_{12}}
\newcommand{\genpriortwoone}{\tilde{q}_{21}}
\newcommand{\priorij}{\beta_i}%{q_{i}}
\newcommand{\priorijst}{q_{i}^*}
\newcommand{\priorjiunder}{\underline{q}_{ji}}
\newcommand{\priorji}{\beta_j}%{q_{j}}
\newcommand{\priorjist}{q_{j}^*}
\newcommand{\priorijunder}{\underline{q}_{ij}}
\newcommand{\fullmprior}{\boldsymbol{\mu}}
\newcommand{\mpriorij}{\mu_{ij}}
\newcommand{\mpriorji}{\mu_{ji}}
\newcommand{\mprioronetwo}{\mu_{12}}
\newcommand{\mpriortwoone}{\mu_{21}}
\newcommand{\fspriorij}[1]{q_{ij}^{#1}}
\newcommand{\fspriorji}[1]{q_{ji}^{#1}}
\newcommand{\fsprioronetwo}[1]{q_{12}^{#1}}
\newcommand{\fspriortwoone}[1]{q_{21}^{#1}}
\newcommand{\fullfsprior}[1]{\fullprior^{#1}}
\newcommand{\priorone}{\beta_1}%{q_1}
\newcommand{\priortwo}{\beta_2}%{q_2}
\newcommand{\prioronetwo}{\beta_1}%{q_{1}}
\newcommand{\priortwoone}{\beta_2}%{q_{2}}
\newcommand{\priorn}{\beta_{n}}%{q_{n}}
\newcommand{\prioronetwounder}{\underline{q}_{12}}
\newcommand{\priortwooneunder}{\underline{q}_{21}}
\newcommand{\fsactmisone}{\overline{\actmisone}}
\newcommand{\fsactmistwo}{\overline{\actmistwo}}
\newcommand{\posti}{Q_i}
\newcommand{\postj}{Q_j}
\newcommand{\postone}{Q_1}
\newcommand{\posttwo}{Q_2}
\newcommand{\belij}{\vec{Q}_{ij}}
\newcommand{\belji}{\vec{Q}_{ji}}
\newcommand{\belonetwo}{\vec{Q}_{12}}
\newcommand{\beltwoone}{\vec{Q}_{21}}
\newcommand{\belone}{\hat{q}_1}
\newcommand{\beltwo}{\hat{q}_2}
\newcommand{\beliefi}{\tau_i}
\newcommand{\beliefj}{\tau_j}
\newcommand{\beliefone}{\tau_1}
\newcommand{\belieftwo}{\tau_2}
\newcommand{\indices}{_{i=1}^{{n}}}
\newcommand{\sactspaceone}{\actspaceone^{\textnormal{s}}}
\newcommand{\sactspacetwo}{\actspacetwo^{\textnormal{s}}}
\newcommand{\sactspacei}{\actspacei^{\textnormal{s}}}
\newcommand{\spayone}{\payone^{\textnormal{s}}}
\newcommand{\spaytwo}{\paytwo^{\textnormal{s}}}
\newcommand{\spayi}{\payi^{\textnormal{s}}}
\newcommand{\spayj}{\payj^{\textnormal{s}}}
\newcommand{\gamesub}{\game^{\textnormal{s}}}
\newcommand{\fullprog}{\textbf{\prog}}
\newcommand{\newfullprog}{\widetilde{\textbf{\prog}}}
\newcommand{\newfullprogi}{{\newfullprog^i}}
\newcommand{\newfullprogmi}{\newfullprog_{-i}}
\newcommand{\fulleq}{\textbf{\prog}^*}
\newcommand{\actfulleq}{\fullact^*}
\newcommand{\genfulleq}{\tilde{\textbf{\prog}}^*}
\newcommand{\altprog}{\tilde{\textbf{\prog}}}
\newcommand{\fullbel}{\vec{\fullpost}}
\newcommand{\fullprior}{\boldsymbol{\beta}}%{\textbf{q}}
\newcommand{\fullpriorone}{\textbf{q}^{1}}
\newcommand{\fullpriorzero}{\textbf{q}^{0}}
\newcommand{\genfullprior}{\tilde{\textbf{q}}}
\newcommand{\fullburn}{\textbf{b}}
\newcommand{\fullcom}{\textbf{c}}
\newcommand{\progpayh}{U_H}
\newcommand{\progpaya}{U_A}
\newcommand{\progpayi}{U_i}
\newcommand{\progpayj}{U_j}
\newcommand{\progpayip}{U_{i'}}
\newcommand{\progpayjp}{U_{j'}}
\newcommand{\fprogi}{f_i}
\newcommand{\fprogj}{f_j}
\newcommand{\fprogone}{f_1}
\newcommand{\fprogtwo}{f_2}
\newcommand{\payh}{u_H}
\newcommand{\paya}{u_A}
\newcommand{\payi}{u_i}
\newcommand{\payip}{u_{i'}}
\newcommand{\payj}{u_j}
\newcommand{\avgpayi}{\overline{u}_i}
\newcommand{\avgpayj}{\overline{u}_j}
\newcommand{\payii}{u_i^i}
\newcommand{\payji}{u_j^i}
\newcommand{\payij}{u_i^j}
\newcommand{\payjj}{u_j^j}
\newcommand{\avgpayii}{\overline{u}_i^i}
\newcommand{\avgpayji}{\overline{u}_j^i}
\newcommand{\avgpayij}{\overline{u}_i^j}
\newcommand{\fullpostsim}{\fullpost^{S}}
\newcommand{\avgpayjj}{\overline{u}_j^j}
\newcommand{\payone}{u_1}
\newcommand{\paytwo}{u_2}
\newcommand{\payn}{u_n}
\newcommand{\progpayone}{U_1}
\newcommand{\progpaytwo}{U_2}
\newcommand{\progone}{\prog_1}
\newcommand{\progtwo}{\prog_2}
\newcommand{\progn}{\prog_n}
\newcommand{\progmn}{\fullprog_{-n}}
\newcommand{\bestone}{b_1}
\newcommand{\besttwo}{b_2}
\newcommand{\besti}{b_i}
\newcommand{\bestj}{b_j}
\newcommand{\eqbestone}{b^*_1}
\newcommand{\eqbesttwo}{b^*_2}
\newcommand{\eqbesti}{b^*_i}
\newcommand{\eqbestj}{b^*_j}
\newcommand{\fulleqbest}{\textbf{b}^*}
\newcommand{\progij}{\prog_i^j}
\newcommand{\progjj}{\prog_j^j}
\newcommand{\progii}{\prog_i^i}
\newcommand{\progji}{\prog_j^i}
\newcommand{\progonetwo}{\prog_1^2}
\newcommand{\progtwotwo}{\prog_2^2}
\newcommand{\progoneone}{\prog_1^1}
\newcommand{\progtwoone}{\prog_2^1}
\newcommand{\actone}{a_1}
\newcommand{\acttwo}{a_2}
\newcommand{\actn}{a_n}
\newcommand{\genactone}{\tilde{a}_1}
\newcommand{\genacttwo}{\tilde{a}_2}
\newcommand{\genacti}{\tilde{a}_i}
\newcommand{\genactj}{\tilde{a}_j}
\newcommand{\response}{\pi}
\newcommand{\fsactone}{\overline{a}_1}
\newcommand{\fsacttwo}{\overline{a}_2}
\newcommand{\fsacti}{\overline{a}_i}
\newcommand{\fsactj}{\overline{a}_j}
\newcommand{\fspoltwo}{\overline{\pi}_2}
\newcommand{\maximize}{M}
\newcommand{\acti}{a_i}
\newcommand{\actj}{a_j}
\newcommand{\payoff}{\mathbf{x}}
\newcommand{\fullbeliefs}{\boldsymbol{\tau}}
\newcommand{\simul}{s}
\newcommand{\simulone}{s_1}
\newcommand{\simultwo}{s_2}
\newcommand{\postij}{Q_{ij}}
\newcommand{\postji}{Q_{ji}}
\newcommand{\ppostij}{q_{ij}^{\partj}}
\newcommand{\ppostji}{q_{ji}^{\parti}}
\newcommand{\ppostonetwo}{q_{12}^{\parttwo}}
\newcommand{\pposttwoone}{q_{21}^{\partone}}
\newcommand{\ppostijone}{q_{ij}^{1}}
\newcommand{\ppostijzero}{q_{ij}^{0}}
\newcommand{\ppostjione}{q_{ji}^{1}}
\newcommand{\ppostjizero}{q_{ji}^{0}}
\newcommand{\ppostonetwozero}{q_{12}^{0}}
\newcommand{\pposttwoonezero}{q_{21}^{0}}
\newcommand{\postijj}{Q_{ij}^j}
\newcommand{\postjii}{Q_{ji}^i}
\newcommand{\fullpost}{\textbf{Q}}
\newcommand{\fullpostonetwo}{\textbf{Q}^{12}}
\newcommand{\fullposttwoone}{\textbf{Q}^{21}}
\newcommand{\fullbelonetwo}{{\textbf{Q}}^{12}}
\newcommand{\fullbeltwoone}{{\textbf{Q}}^{21}}
\newcommand{\fullbelsim}{{\textbf{Q}}^{S}}
\newcommand{\postonetwo}{Q_{12}}
\newcommand{\posttwoone}{Q_{21}}
\newcommand{\postonetwotwo}{Q_{12}^2}
\newcommand{\posttwooneone}{Q_{21}^1}
\newcommand{\gpostonetwo}{\tilde{Q}_{12}}
\newcommand{\gposttwoone}{\tilde{Q}_{21}}
\newcommand{\gpostij}{\tilde{Q}_{ij}}
\newcommand{\gpostji}{\tilde{Q}_{ji}}
\newcommand{\simuli}{s_i}
\newcommand{\simulj}{s_j}
\newcommand{\bri}{\pi^*_i}
\newcommand{\brj}{\pi^*_j}
\newcommand{\brone}{\pi^*_1}
\newcommand{\brtwo}{\pi^*_2}
\newcommand{\altdefprog}[1]{\underline{#1}}
\newcommand{\willreneg}{\mathcal{R}}
\newcommand{\prog}{p}
\newcommand{\weakpoint}{a}
\newcommand{\attackguess}{b}
\newcommand{\revelation}{\mathcal{R}}
\newcommand{\partialrevelation}{\mathcal{T}}
\newcommand{\eqsolver}{E}
\newcommand{\nullout}{\emptyset}
\newcommand{\penaltybehav}{c}
\newcommand{\penaltyrat}{d_1}
\newcommand{\penaltyeq}{d_2}
\newcommand{\penaltypref}{d_3}
\newcommand{\one}{^{(1)}}
\newcommand{\two}{^{(2)}}
\newcommand{\ksup}{^{(k)}}
\newcommand{\xk}{\mathbf{x}\ksup}
\newcommand{\xkap}{\mathbf{x}(\kappa)}
\newcommand{\xkapk}{\mathbf{x}(\kappa\ksup)}
\newcommand{\ykap}{\mathbf{y}(\kappa)}
\newcommand{\ykapk}{\mathbf{y}(\kappa\ksup)}
\newcommand{\act}{a}
\newcommand{\br}{\text{BR}}
\newcommand{\fair}{F}
\newcommand{\dthrfunc}{T_D}
\newcommand{\sthrfunc}{T_S}
\newcommand{\barg}{B}
\newcommand{\funci}{g_i}
\newcommand{\funcj}{g_j}
\newcommand{\funcone}{g_1}
\newcommand{\functwo}{g_2}
\newcommand{\csprog}{\mathbf{csr}}
\newcommand{\actset}{A}
\newcommand{\agsetpmm}[0]{PMM}
\newcommand{\A}[0]{$A$}
\newcommand{\B}[0]{$B$}



\usepackage{makecell}
\usepackage{algorithm}
\usepackage[noend]{algpseudocode}%http://ctan.org/pkg/algorithmicx

\usepackage{tikz}

\usetikzlibrary{bayesnet, arrows.meta, shapes, positioning,
patterns, patterns.meta}
\usetikzlibrary{matrix, fit}
\usetikzlibrary{decorations.markings}
\tikzset{
    myarrow/.style={
        -Stealth,
        shorten >=8pt,
        shorten <=8pt,
    },
    mymatrix/.style={
        matrix of math nodes
    },
}


\makeatletter
\newcommand{\textlabel}[2]{%
  \@bsphack
  \csname phantomsection\endcsname % in case hyperref is used
  \def\@currentlabel{#1}{\label{#2}}%
  \@esphack
}
\makeatother

\newtheorem{theorem}{Theorem}
\newtheorem{lemma}{Lemma}
\newtheorem{proposition}[theorem]{Proposition}
\newtheorem{corollary}[theorem]{Corollary}
\newtheorem{fact}[theorem]{Fact}
\theoremstyle{definition}
\newtheorem{exmp}{Example}[section]
\newtheorem{principle}[theorem]{Principle}
\newenvironment{sketch}{\paragraph{\normalfont \textsc{Proof Sketch.}}}{\hfill$\square$ \\}

\usepackage{mathtools}
\newtheorem{remark}[theorem]{Remark}
% --- end inlined defaults.sty ---


\usepackage{balance} % for balancing columns on the final page



\setcopyright{ifaamas}
\acmConference[AAMAS '25]{Proc.\@ of the 24th International Conference
on Autonomous Agents and Multiagent Systems (AAMAS 2025)}{May 19 -- 23, 2025}
{Detroit, Michigan, USA}{A.~El~Fallah~Seghrouchni, Y.~Vorobeychik, S.~Das, A.~Nowe (eds.)}
\copyrightyear{2025}
\acmYear{2025}
\acmDOI{}
\acmPrice{}
\acmISBN{}




\acmSubmissionID{120}


\title[AAMAS-2025 Formatting Instructions]{Safe Pareto Improvements
for \\
Expected Utility
Maximizers in Program Games}


\author{Anthony DiGiovanni}
\affiliation{
  \institution{Center on Long-Term Risk}
  \city{London}
  \country{United Kingdom}}
\email{anthony.digiovanni@longtermrisk.org}

\author{Jesse Clifton}
\affiliation{
  \institution{Center on Long-Term Risk}
  \city{London}
  \country{United Kingdom}}
\email{jesse.clifton@longtermrisk.org}

\author{Nicolas Macé}
\affiliation{
  \institution{Center on Long-Term Risk}
  \city{London}
  \country{United Kingdom}}
\email{nicolas.mace@longtermrisk.org}



\begin{abstract}
Agents in mixed-motive coordination problems such as Chicken may fail to coordinate on a Pareto-efficient outcome. 
Safe Pareto improvements (SPIs) were originally proposed 
to mitigate miscoordination in cases where 
players lack probabilistic beliefs as to how their 
agents will play a game;
agents are instructed to behave so as to guarantee a Pareto improvement on how they would play by default.
More generally, SPIs may be defined as transformations of strategy profiles such that all players are necessarily better off under the transformed profile.
In this work, we investigate
the extent to which SPIs can reduce downsides of miscoordination
between expected utility-maximizing agents. We 
consider games in which players submit computer programs
that can condition their decisions on each other's code, 
and use this property to construct SPIs using programs
capable of \textit{renegotiation}. 
We first
show that under mild conditions on players’ beliefs,
each player always prefers to use renegotiation.
Next, we show that
under similar assumptions,
each player always prefers
to be willing to renegotiate  
at least to the point at which 
they receive the lowest
payoff they can attain in any efficient
outcome.
Thus 
subjectively optimal play
guarantees players at least these payoffs,
without the need for coordination
on specific Pareto improvements.
\end{abstract}




\keywords{program equilibrium; bargaining; Pareto efficiency; Cooperative AI}


         
\newcommand{\BibTeX}{\rm B\kern-.05em{\sc i\kern-.025em b}\kern-.08em\TeX}


\begin{document}


\pagestyle{fancy}
\fancyhead{}


\subsection{Revised Formal Definition of
CSR}
\label{sec:sub:components}

Here 
is the formal definition of CSR
programs, with the PMP-renegotiation
variant built into the same program class.
For a
set-valued
renegotiation
function~$\renegi$,
we define
the space of ordinary
CSR
programs
$\fsspaceireni$
as before: these are programs
with the structure
of
Algorithm~\ref{alg:csr},
for some default~$\defprogi$,
that use~$\renegi$ as their base
renegotiation function.

We also jointly define a distinguished subclass
$\prnspacei(\defrenegi)$,
called
\termemph{PMP-renegotiation (PRN) programs}.
A PRN program has the same kind of default
program as an ordinary CSR program, but it also has
a \termemph{reference} set-valued renegotiation function
$\defrenegi$.
Against CSR programs outside~$\prnspacej$,
it is evaluated exactly as the ordinary CSR program
with reference function~$\defrenegi$.
Against programs inside~$\prnspacej$,
it is evaluated as the PMP-extension of~$\defrenegi$.
This is the structural constraint we need:
adding the PMP is conditional on the counterpart also using a PRN program.
Thus a counterpart who switches to a non-PRN CSR program in order to
``punish'' the PMP-extension only faces the reference CSR program,
not the PMP-extension.

Let~$\fullsvspace$ be the set of all
set-valued renegotiation 
functions.
For each~$i$, let
$\fullfsspacei = \bigcup_{\renegi \in \fullsvspace} \fsspaceireni$
and
$\prnspacei = \bigcup_{\defrenegi \in \fullsvspace}\prnspacei(\defrenegi)$.
Let
$\allprogreneg = \fullfsspaceone \times \fullfsspacetwo$
and
$\allprogprn = \prnspaceone \times \prnspacetwo$.
The \selfunction{} function~$\renegb$
is given to
the players
(and we suppress dependence of~$\fullfsspacei$ and~$\prnspacei$ on~$\renegb$ for simplicity).
For either an ordinary CSR program or a PRN program~$p$,
write~$\bar{\reneg}^{p}$ for its base/reference function:
for an ordinary CSR program this is its set-valued renegotiation function,
and for a PRN program this is its reference set-valued renegotiation function.
For readability, Algorithm~\ref{alg:csr} writes these two functions as $\bar{\reneg}^{1}$ and $\bar{\reneg}^{2}$.

\begin{algorithm}
\caption{
Conditional set-valued
renegotiation / PRN
program $\progi$, for some default~$\defprogi$}
\label{alg:csr}
\begin{algorithmic}[1]
\Require Counterpart program $\progj$
\If{$\progj$ is an ordinary CSR or PRN program, with base/reference function $\bar{\reneg}^{j}$}
    \State $\deffullact \leftarrow \fullact(\deffullprog)$
    \If{$\progi \in \prnspacei$ and $\progj \in \prnspacej$}
        \State ${I}^{\mathrm{def}} \leftarrow \bar{\reneg}^{1}(\bar{\reneg}^{2}, \deffullact) \cap \bar{\reneg}^{2}(\bar{\reneg}^{1}, \deffullact)$
        \If{${I}^{\mathrm{def}} \neq \emptyset$}
            \State $\fullact^{\mathrm{def}} \leftarrow \renegb({I}^{\mathrm{def}})$
        \Else
            \State $\fullact^{\mathrm{def}} \leftarrow \deffullact$
        \EndIf
        \State ${I} \leftarrow \left(\bar{\reneg}^{1}(\bar{\reneg}^{2}, \deffullact) \cup \yudprojone(\fullact^{\mathrm{def}})\right)$
        \Statex \hspace{1.7cm}$\cap \left(\bar{\reneg}^{2}(\bar{\reneg}^{1}, \deffullact) \cup \yudprojtwo(\fullact^{\mathrm{def}})\right)$ \Comment{PMP-extended agreement set}
    \Else
        \State ${I} \leftarrow \bar{\reneg}^{1}(\bar{\reneg}^{2}, \deffullact) \cap \bar{\reneg}^{2}(\bar{\reneg}^{1}, \deffullact)$ \Comment{Reference CSR agreement set}
    \EndIf
    \If{${I} \neq \emptyset$}
        \State $\deffullact \leftarrow \renegb({I})$ \Comment{Renegotiation outcome}
    \EndIf
    \State \Return $\defacti$ \Comment{Play renegotiation outcome, or default}
\Else
     \State \Return $\defprogi(\progj)$  % \Comment{Play default}
\EndIf
\end{algorithmic}
\end{algorithm}


\subsection{Guaranteeing PMM Payoffs Using
CSR}
\label{sec:sub:pmm}


Suppose
players always include
more outcomes in their renegotiation
sets if their expected
utility is unchanged.
So in particular, to show that in subjective equilibrium players use programs
that guarantee at least the PMM,
it will suffice to show that
they weakly prefer these programs.

Then, 
we will show in
Theorem~\ref{prop:yudpoint}
that
under
mild assumptions on players' beliefs,
for any program
that
does not guarantee
a player at least
their PMM payoff,
there is a corresponding PRN
program the player prefers
that \textit{does}
guarantee their PMM
payoff.
We prove this
result
by constructing programs that are identical
to the programs
players would otherwise use
against non-PRN counterparts,
but that add the player's PMP
against PRN counterparts.
For a program $\progi$, we call this
modified program the
\termemph{PMP-CSR extension} of~$\progi$
(Definition~\ref{def:pmp-extension}).

This result requires only the same kind of
no-punishment assumption used for ordinary renegotiation
(Assumption~\ref{assum:nopunish}),
applied to CSR and PRN programs:
programs that do not use the CSR/PRN machinery are assumed not to
punish such programs when no renegotiation occurs.
The PRN construction removes the need for a separate assumption that
counterparts do not punish PMP-extension itself.
By construction, if the counterpart is outside~$\prnspacej$,
player~$i$'s PRN program is evaluated exactly as its reference CSR program.
If the counterpart is inside~$\prnspacej$, then the counterpart is also using
the PMP-extension branch.
Thus the PMP points are added only in the case where the structural response
is symmetric, rather than in cases where a counterpart has moved to a
more restrictive non-PRN renegotiation set.


For Theorem~\ref{prop:yudpoint}
we also assume the \selfunction{} function is transitive.
This is an intuitive property:\ If outcomes are added to the agreement
set that make \textit{all} players weakly better off
than the default renegotiation outcome,
the new renegotiation outcome should be weakly better for all players.

\par The remainder of this subsection provides the formal 
details for the statement of Theorem \ref{prop:yudpoint}, 
and a sketch of the proof.


\multilinecomment{
\begin{figure}
    \centering

  \begin{tabular}{cc}

    \begin{tikzpicture}[scale=1]
    \usetikzlibrary {shapes.geometric} 

\draw[thick,<->] (-0.4, 0) -- (2.8, 0) node[right] {$x$};
\draw[thick,<->] (0, -0.4) -- (0, 2.8) node[right] {$y$};

\draw[gray, fill=gray] (0, 0) circle (0.08);
\draw[gray, fill=gray] (2.4, 0.8) circle (0.08);
\draw[gray, fill=gray] (0.8, 2.4) circle (0.08);
\draw[gray, fill=gray] (0.8, 0.8) circle (0.08);

\draw[gray, thick] (0,0) -- (2.4,0.8) -- (0.8,2.4) -- cycle;

\draw[thick,dashed,-] (0.8, 0.8) -- (2.4, 0.8);

\fill[pattern=north east lines, pattern color=gray] 

  (0,0) -- (0.15, 0.45) -- (1.35,0.45) -- cycle;
  
\fill[pattern=north west lines, pattern color=gray]
    (0,0) -- (\pointx, \tritop) -- (\pointx,\tribottom) -- cycle;

   \draw[black, fill=black] (\pointx, 0.8) circle (0.06);

\end{tikzpicture}



    &




        \begin{tikzpicture}[scale=1]
    \usetikzlibrary {shapes.geometric} 

\draw[thick,<->] (-0.4, 0) -- (2.8, 0) node[right] {$x$};
\draw[thick,<->] (0, -0.4) -- (0, 2.8) node[right] {$y$};

\draw[gray, fill=gray] (0, 0) circle (0.08);
\draw[gray, fill=gray] (2.4, 0.8) circle (0.08);
\draw[gray, fill=gray] (0.8, 2.4) circle (0.08);
\draw[gray, fill=gray] (0.8, 0.8) circle (0.08);

\draw[gray, thick] (0,0) -- (2.4,0.8) -- (0.8,2.4) -- cycle;

\draw[thick,dashed,-] (0.8, 0.8) -- (2.4, 0.8);

\fill[pattern=north east lines, pattern color=gray] 

  (0,0) -- (0.15, 0.45) -- (1.35,0.45) -- cycle;
  
\fill[pattern=north west lines, pattern color=gray]
    (0,0) -- (\pointx-0.2, \tritop-0.55) -- (\pointx-0.2,\tribottom-0.08) -- cycle;

   \draw[black, fill=black] (\pointx-0.2, 0.45) circle (0.06);


\node[anchor=north east] at (4,3.5) {
\fbox{
\begin{tabular}{ll}
    \tikz\fill[pattern=north east lines,pattern color=gray] (0,0) rectangle (0.25,0.25); & P1's %~$\renegoneone$ 
    set \\
 \tikz\fill[pattern=north west lines,pattern color=gray] (0,0) rectangle (0.25,0.25); & P2's set \\
 \tikz \node [star,
 star point height=.01cm,
 minimum size=0.01cm, 
 star point ratio=0.3,
 draw]
       at (\pointx,\pointx) {}; & PMM \\
\end{tabular}
}
};

  
\end{tikzpicture}
  

\\
    \end{tabular}
    
    \caption{Caption}
    \label{fig:enter-label}
\end{figure}
}


\begin{definition}\label{def:pmp-extension}
For any ordinary CSR program~$\progi \in \fsspaceireni$,
the \textbf{PMP-CSR extension}~$\widetilde{\progi} \in \prnspacei(\renegi)$
is the PRN program with the same default program as~$\progi$
and reference function~$\renegi$.
Thus, if the counterpart is not in~$\prnspacej$,
$\widetilde{\progi}$ is evaluated exactly as~$\progi$.
If the counterpart is a PRN program in~$\prnspacej(\defrenegj)$,
write~$\deffullact$ for the action profile induced by the two defaults,
let
\[
    {I}^{\mathrm{def}}
    = \renegi(\defrenegj, \deffullact)
    \cap \defrenegj(\renegi, \deffullact),
\]
and let~$\fullact^{\mathrm{def}} = \renegb({I}^{\mathrm{def}})$ if
${I}^{\mathrm{def}} \neq \emptyset$, and
$\fullact^{\mathrm{def}}=\deffullact$ otherwise.
Then the active renegotiation set used by~$\widetilde{\progi}$ is
\[
    \renegi(\defrenegj, \deffullact)
    \cup \yudproji(\fullact^{\mathrm{def}}).
\]
\end{definition}

\begin{assumption}
     We say that 
     players 
     with beliefs $\fullprior$ 
     \textbf{are certain that CSR/PRN won't be punished} 
     if the following holds:
     \begin{enumerate}[(i)]
      \item Suppose either $\progi$ is in a subjective equilibrium
     of~$\game(\allprog \cup \allprogreneg \cup \allprogprn)$,
     or $\progi$ is in the support of $\priorji$.
      Suppose $\progi \notin \fsspacei \cup \prnspacei$. 
         Then for any $\progj \in  \fsspacej \cup \prnspacej$,
         we have 
         $\progi(\progj) = \progi(\defprogj)$.
         \label{subassum:csrwlog}
     \end{enumerate}
     \label{assum:nopunish}
\end{assumption}

\begin{figure}
    \centering
     \begin{tikzpicture}[scale=1]
    \usetikzlibrary {shapes.geometric} 

\draw[thick,<->] (-0.4, 0) -- (2.8, 0) node[right] {$x$};
\draw[thick,<->] (0, -0.4) -- (0, 2.8) node[right] {$y$};

\draw[gray, fill=gray] (0, 0) circle (0.08);
\draw[gray, fill=gray] (2.4, 0.8) circle (0.08);
\draw[gray, fill=gray] (0.8, 2.4) circle (0.08);
\draw[gray, fill=gray] (0.8, 0.8) circle (0.08);
\draw[gray, thick] (0,0) -- (2.4,0.8) -- (0.8,2.4) -- cycle;

\fill[pattern=north east lines, pattern color=gray] 
(0,0) -- (0.4, 1.2) -- (0.4, 0.6) -- (1.8,0.6) -- cycle;
  
  
\fill[pattern=north west lines, pattern color=gray]
   (0,0) -- (0.8, 2.4) -- (0.8,0.3) -- cycle;

\draw[black,fill=black] (-0.08, -0.08) rectangle (0.08, 0.08);

 \draw[black, fill=black] (0.4, 1.2) circle (0.06);

\draw[thick,dashed,-] (0.8, 1.2) -- (2, 1.2);


\node[anchor=north east] at (7,3) {
\fbox{
\begin{tabular}{ll}
 \tikz\node[rectangle,draw,fill=black] at (0.3,0.3) {}; & $\fullact(\deffullprog)$ \\
    \tikz\fill[pattern=north east lines,pattern color=gray] (0,0) rectangle (0.3,0.3); & $\renegone(\renegtwo, \fullact(\deffullprog))$ \\
 \tikz\fill[pattern=north west lines,pattern color=gray] (0,0) rectangle (0.3,0.3); & $\renegtwo(\renegone, \fullact(\deffullprog))$ \\
    \tikz\node[circle,draw,fill=black,scale=0.6] at (0,0) {}; & $\fullact(\fullprog)$ \\ 
 \tikz[baseline]{\draw[dashed] (0,0.1) -- (0.35,0.1);}
 & $\yudprojone(\fullact(\fullprog))$ \\
\end{tabular}
}
};

  
\end{tikzpicture}
    \caption{Illustration of the argument for Theorem~\ref{prop:yudpoint}.
    The black circle is the reference CSR renegotiation outcome,
    $\fullact^{\mathrm{def}}$.
    If both players use PRN programs, player~1 adds the black striped
    segment $\yudprojone(\fullact^{\mathrm{def}})$ to their reference
    renegotiation set, and player~2 symmetrically uses their PRN branch.
    If player~2 is not in~$\prnspacetwo$, player~1 instead uses the
    unextended reference set $\renegone(\renegtwo, \fullact(\deffullprog))$.
    Thus a non-PRN counterpart cannot make themselves better off by
    switching to a restrictive renegotiation set in response to PMP points:
    against that counterpart, those PMP points are not added.
    }
    \label{fig:pmp}
    \Description{The default renegotiation outcome a(p) is marked with a black circle. A black striped segment shows Player 1's PMP of this outcome, PMP₁(a(p)), which includes points that make both players weakly better off than a(p) while not making Player 2 strictly better off. The shaded region represents Player 1's renegotiation set RN¹(RN², a(p^def)).}
\end{figure}
\begin{restatable}{theorem}{pmmguarantee}
\label{prop:yudpoint}
Let~$\game(\allprog)$ be a program game,
and~$\renegb$ be any transitive \selfunction{} function.
Suppose
the action sets of~$\game$
are continuous, so
that for any~$\fullact \in \allact$,
player~$i$'s PMP of that action profile
$\yudproji(\fullact)$
is nonempty.
Let~$\fullprior$ be any belief profile
satisfying the assumption 
that players are certain that CSR/PRN won't be punished
(Assumption~\ref{assum:nopunish}).

 Then,
 for any subjective equilibrium~$(\fullprog, \fullprior)$ of
 $\game(\allprog \cup \allprogreneg \cup \allprogprn)$
 where $\progpayi(\fullprog) < \pmmi$ 
 for some~$i$,
 there exists~$\fullprog' \in \allprogprn$
 such that:
    \begin{enumerate}
        \item For all~$i$, $\progi'$ is a PMP-CSR extension
        with the same default behavior as~$\progi$.
    \item $\fullprogpay(\fullprog') \succeq \pmm$.
    \item $(\fullprog', \fullprior)$
      is a subjective equilibrium
      of  $\game(\allprog \cup \allprogreneg \cup \allprogprn)$.
    \end{enumerate}
\end{restatable}
\begin{sketch}
Assumption~\ref{assum:nopunish}
implies that any program outside the CSR/PRN class is payoff-equivalent,
against the relevant beliefs,
to an ordinary CSR program that has the original program as its default
and always returns the empty renegotiation set.
So it suffices to consider ordinary CSR programs and PRN programs.

Take any ordinary CSR program~$\progi$ with base function~$\renegi$,
and replace it with its PMP-CSR extension~$\newprogi \in \prnspacei(\renegi)$.
Against any counterpart outside~$\prnspacej$,
Algorithm~\ref{alg:csr} evaluates~$\newprogi$ exactly as the reference
CSR program~$\progi$.
Against a counterpart inside~$\prnspacej$,
both programs use their PRN branches: each starts from the reference CSR
agreement set and adds their PMP of the reference CSR outcome.
The only new outcomes that can become relevant for player~$i$ are therefore
weak Pareto improvements on the reference CSR outcome.
Because the \selfunction{} function is transitive,
player~$i$ is weakly better off in this case as well.
Thus each player weakly prefers their PMP-CSR extension.

When all players use PMP-CSR extensions,
the agreement set contains each player's PMP of the reference CSR outcome.
Any agreement-set outcome giving some player less than their PMM payoff is
Pareto-dominated by an outcome in the intersection of these PMPs, and so
cannot be selected by a \selfunction{} function that chooses an efficient
point from the agreement set.
Therefore the resulting program profile guarantees every player at least
PMM payoff, and the weak-improvement argument above preserves subjective
equilibrium.
\end{sketch}


\textit{Remark:} Notice that the argument above does not require players
to refrain from using programs that implement other kinds of SPIs.
A player might, for example, use a CSR program whose reference set is
self-favoring rather than PMP-based.
The theorem says that, whatever reference CSR behavior the player starts from,
they weakly prefer the corresponding PRN wrapper:
outside the PRN class the wrapper behaves like the reference program, and
inside the PRN class it adds the PMP.
This also explains why the PMP is the natural extension rather than a larger
set of Pareto improvements.
Adding outcomes that are better for the counterpart and worse for the focal
player would no longer be protected by the same structural argument: against
PRN counterparts those outcomes might be selected, and against non-PRN
counterparts the PRN wrapper deliberately does not add them.


\

We can now formalize the claim that
PRN-augmented CSR
is an SPI
that partially solves
SPI selection:
for any reference profile of set-valued renegotiation functions,
the mapping from programs~$\fullprog$
to PRN programs using~$\fullprog$ as defaults
is an SPI that guarantees players
their PMM payoffs under the assumptions of
Theorem~\ref{prop:yudpoint}.

\begin{proposition}
\label{prop:itspi}
For $i=1,2$,
for some \selfunction{} function~$\renegb$,
define 
$\fsprogi^{\renegi}: \progspacei \rightarrow \prnspacei(\renegi)$ such that,
for each $\progi \in \progspacei$, $\fsprogi^{\renegi}(\progi)$
is the PMP-CSR extension with reference function~$\renegi$ and
$\defprog{\fsprogi^{\renegi}(\progi)} = \progi$.
Then, 
under the assumptions of 
Theorem~\ref{prop:yudpoint},
for any~$(\renegone, \renegtwo)$:
\begin{enumerate}
    \item The function
    $\fullfs^{\reneg}: \fullprog \mapsto (\fsprogone^{\renegone}(\progone), \fsprogtwo^{\renegtwo}(\progtwo))$ is an SPI.
    \item For all~$i$,
    $\progpayi(\fullfs^{\reneg}(\deffullprog)) \geq \max\{\progpayi(\deffullprog), \pmmi\}$.
\end{enumerate}
\end{proposition}

\begin{proof}
This follows
immediately
from the argument used to prove 
Theorem~\ref{prop:yudpoint}.
\end{proof}


\section{Proof of Theorem~\ref{prop:yudpoint} }

Denote 
the \textbf{agreement set} as 
     $\intsect(\renegi, \renegmi, \fullact) = \bigcap_{j=1}^{n} \renegj(\renegmj, \fullact)$.

For the ${n}$-player proof, the PRN branch of
Algorithm~\ref{alg:csr} is read as follows.
For a profile of base/reference functions~$(\bar{\reneg}^{j})_{j=1}^{n}$
and default action profile~$\deffullact$, let
\[
    {I}^{\mathrm{def}}
    = \bigcap_{j=1}^{n} \bar{\reneg}^{j}(\bar{\reneg}^{-j}, \deffullact),
\]
and let~$\fullact^{\mathrm{def}}=\renegb({I}^{\mathrm{def}})$ if
${I}^{\mathrm{def}}\neq\emptyset$, and
$\fullact^{\mathrm{def}}=\deffullact$ otherwise.
When all players are in their PRN classes, player~$i$'s active
renegotiation set is
\[
    \bar{\reneg}^{i}(\bar{\reneg}^{-i}, \deffullact)
    \cup \yudproji(\fullact^{\mathrm{def}}).
\]
When at least one player is outside the PRN class, the algorithm uses the
reference CSR sets~$\bar{\reneg}^{i}(\bar{\reneg}^{-i},\deffullact)$
without adding PMP points.

\pmmguarantee*

\begin{proof}
Let~$(\fullprog, \fullprior)$
be a subjective equilibrium
of $\game(\allprog \cup \allprogreneg \cup \allprogprn)$.
First, any program outside the ordinary CSR and PRN classes can be replaced,
without changing its payoff against the relevant beliefs,
by an ordinary CSR program whose default is the original program and whose
renegotiation set is always empty.
Indeed, if the counterpart profile does not enter the CSR/PRN branch, the new
program directly calls its default.
If the counterpart profile does enter that branch, the empty agreement set
again makes the new program play its default; by
Assumption~\ref{assum:nopunish}, the counterpart programs also respond as
they would to the default.
Thus we may restrict attention to ordinary CSR programs and PRN programs.

For each player~$i$, let~$\progi^0$ be the reference ordinary CSR program
associated with~$\progi$: if~$\progi$ is ordinary CSR, let
$\progi^0=\progi$; if~$\progi$ is PRN, let~$\progi^0$ be the ordinary CSR
program with the same default and reference function.
Let~$\progi'$ be the PMP-CSR extension of~$\progi^0$.
We show that~$\progi'$ is weakly better for player~$i$ against every
counterpart profile in the support of~$\priorij$.

If the counterpart profile is not entirely composed of PRN programs, then by
construction Algorithm~\ref{alg:csr} evaluates~$\progi'$ exactly as the
reference CSR program~$\progi^0$.
So the outcome is unchanged relative to the reference behavior.
If the counterpart profile is entirely composed of PRN programs, then every
player's active set is their reference CSR set plus their PMP of the reference
CSR outcome~$\fullact^{\mathrm{def}}$.
The only outcomes added to the agreement set that can change player~$i$'s
payoff are weak Pareto improvements on~$\fullact^{\mathrm{def}}$.
Since~$\renegb$ is transitive, selecting from the enlarged agreement set makes
player~$i$ weakly better off than selecting from the reference agreement set.
Thus~$\progi'$ is weakly better for player~$i$ against every program profile in
the support of their beliefs.
Applying this argument for each player, the profile
$\fullprog'=(\progi')_{i=1}^{n}$ is also a subjective equilibrium.

It remains to show that~$\fullprog'$ guarantees PMM payoffs.
Let~$\fullact^{\mathrm{def}}$ be the reference CSR outcome for the defaults
and reference functions of~$\fullprog'$.
Since all players are PRN at~$\fullprog'$, the agreement set contains
\[
    \bigcap_{i=1}^{n} \yudproji(\fullact^{\mathrm{def}}).
\]
Every point in this intersection gives every player at least
$\max\{\payi(\fullact^{\mathrm{def}}),\pmmi\}$.
Therefore any point in the agreement set that gives some player less than
PMM is Pareto-dominated by a point in this intersection.
Because~$\renegb$ selects a point that is Pareto-efficient within the
agreement set, it cannot select such a below-PMM point.
Hence~$\fullprogpay(\fullprog') \succeq \pmm$, as required.
\end{proof}





\end{document}


